\section{Introduction}

Digital twins connect observations of a production system with executable models that support monitoring and operational decisions. Kritzinger et al.~\cite{Kritzinger2018} distinguish digital models, digital shadows, and digital twins by their degree of data integration. This distinction matters for conveyor-based handling: a geometric model can reproduce equipment layout, but supervisory decisions additionally require a representation of workpiece motion, sensor events, and the information available to the controller. The application considered here comprises two conveyors, a pneumatic transfer pusher, an XY gantry with a vertical actuator, and a vacuum gripper. Its operation motivates a specific question: when a photoelectric observation is missing or delayed, should transport continue for another control interval, or should the supervisor request a hold?

\subsection{Related Research}

\paragraph{Digital twins for fault diagnosis.}
Ademujimi and Prabhu~\cite{Ademujimi2022} combined equipment simulation with discrete-event simulation to generate fault interventions for learning Bayesian network structures. They evaluated the approach using a laboratory assembly cell containing four industrial robots, demonstrating how a twin can support diagnostic model development when fault data are difficult to obtain. Their reported future work includes evaluating sensitivity to deviations in twin parameters and extending validation to a manufacturing plant. This establishes both a useful precedent for generating fault scenarios and a relevant development requirement: diagnostic behavior learned from a twin must be assessed against model mismatch. For the present handling cell, a further task is to connect uncertain observations to an immediate operational decision; recovering a diagnostic graph and deciding whether a moving workpiece may advance are different evaluation objectives.

\paragraph{Digital twins for disturbance handling and reconfiguration.}
Fu et al.~\cite{Fu2025} integrated process monitoring, capability descriptions, optimization, and automatically generated simulation to redistribute manufacturing operations under disturbances. Their battery-line case study uses measured production data, while the performance simulator represents operation durations and material flow with physical movements abstracted. It demonstrates reconfiguration at the level of task assignment and line performance. Importantly for the present study, the case treats data-collection delays as negligible relative to operation durations. The authors also identify spatial constraints and experimental testbed validation as future work. A conveyor cell presents a complementary question: observation delays must instead be evaluated relative to the workpiece's remaining travel distance. Reassigning operations is therefore a useful precedent for adaptation, but it does not by itself specify when uncertain transport should be interrupted.

\paragraph{Network-aware twins and information freshness.}
Bellavista et al.~\cite{Bellavista2021} proposed Application-driven Digital Twin Networking middleware to support heterogeneous industrial communication and dynamic management of edge-network resources through software-defined networking. Their work establishes that network awareness is already part of industrial twin research; it cannot be claimed as a new concept here. At the control level, Hespanha et al.~\cite{Hespanha2007} review the effects of sampling, communication delays, and packet losses, while Yates et al.~\cite{Yates2021} explain information freshness through the age of timestamped updates. These results motivate explicit treatment of observation timing. The additional application requirement is to translate information age into motion uncertainty and a supervisory action: a recently received message may still describe an old workpiece position, and a missing transition alone cannot distinguish a failed sensor from delayed delivery or unexpected motion.

\subsection{Research Gap and Study Scope}

The studies discussed above provide complementary foundations for diagnostic learning, production reconfiguration, and network-aware information management. The engineering problem addressed here lies at their interface: selecting an admissible next action in a sparsely observed handling cell when sensing and communication are both imperfect. This is a focused synthesis of selected relevant studies, not evidence that the complete literature lacks such integration. The present research question is whether timestamped observations and bounded motion predictions can support a continue-or-hold decision without access to the true workpiece state inside the simulator. Its evaluation must distinguish delayed information from physical anomalies, enforce a finite computation budget, and compare decisions under matched operating and network conditions.

The implemented framework separates plant simulation, sensing, communication, estimation, sequence control, prediction, and evaluation. MuJoCo provides the multibody and contact simulation environment~\cite{Todorov2012}. Received photoelectric events anchor a one-workpiece position estimator after passing through configurable delay, jitter, and packet loss. Ground truth is reserved for observation generation and evaluation. A reduced one-dimensional predictor evaluates stopping trajectories with position and velocity offsets, reaction delays, and alternative continuation or hold requests. A computation-budget overrun produces a hold request. This model is an engineering approximation awaiting calibration; it is neither a full-contact rollout for every candidate nor a certified physical safety guarantee.

The study contributes an executable architecture with explicit separation of controller information from evaluation state; a predictive supervisory rule that incorporates observation age and bounded motion uncertainty under a computation budget; and paired experiments with recorded configurations and a shared network seed. Four acceptance cases, comprising seven simulation runs, provide initial descriptive evidence. Both nominal runs complete placement in 15.46~s of simulated time. In the tested photoelectric sensor fault, predictive supervision requests a hold at 6.72~s, compared with 6.76~s for the sequence-control timeout. These observations establish an implemented decision path, rather than general fault-detection accuracy or a cycle-time improvement.

Further development must connect this initial supervision to validated software reconfiguration. Priorities are to calibrate geometry and stopping dynamics, quantify the ability to distinguish sensor faults from communication disturbances, and evaluate recovery policies such as sensor substitution or conveyor-speed adjustment. Such policies require explicit conditions for entering and leaving degraded operation, followed by repeated experiments and live programmable logic controller integration. The evaluated plant still uses an uncalibrated baseline geometry; subsequent CAD measurements have not yet been propagated into those experiments. Automatic recovery and continued production after reconfiguration therefore remain research objectives, while continuation and a latched hold are the capabilities demonstrated at this stage.
